\section{Karten - AI Evaluierung}
Siehe \href{sec:KI}{KI Nutzung} für details, Link zum\href{https://gemini.google.com/share/b40061c4f3a6}{Prompt}\\

\noindent
Hallo! Als potenzieller Nutzer finde ich die grundlegende Idee wirklich spannend. Abstrakte Datensätze (wie Baumkataster oder Luftmessstellen) geografisch direkt in meiner eigenen Nachbarschaft entdeckbar zu machen, ist ein toller Ansatz, um städtische Daten für den Alltag greifbar zu machen.
\\

\noindent
Hier ist mein konstruktives Feedback zu eurem Konzept und dem Prototyp:
\\

\noindent
    \emph{Übersichtlichkeit der Karte (Startseite):} Das Design wirkt durch Material 3 schön clean. Ich frage mich jedoch, wie das System skaliert. Wenn es dutzende Datensätze gibt, wird die Karte schnell mit Pins überladen sein. Es braucht eine sehr intuitive Filterfunktion oder Kategorien (z. B. Umwelt, Verkehr, Historisches), damit man gezielt suchen kann und nicht von Icons erschlagen wird.
\\

\noindent
    \emph{Der Medienbruch (Mein größter Kritikpunkt):} In Abschnitt 6.2 schreibt ihr, dass die Datensätze in externen Apps (z. B. einem PDF-Leser) geöffnet werden. Das unterbricht das Nutzererlebnis massiv. Wenn ich auf der Karte einen Pin für eine Luftmessstelle anklicke, erwarte ich, die Werte direkt in eurer App zu sehen. Die App sollte ein Visualisierungstool sein und nicht nur ein Dateimanager, der mich aus der App wirft.
\\

\noindent
    \emph{Wording der Navigation:} Der mittlere Screen im Prototyp heißt "Downloads", wirkt aber eher wie ein Katalog aller verfügbaren Datensätze, aus dem man Dinge erst herunterlädt. Ein Begriff wie "Datenkatalog", "Entdecken" oder "Datensätze" wäre hier intuitiver.
\\

\noindent
\emph{Fazit:} Ein optisch sehr gelungener Prototyp mit einer starken Grundidee. Um echten Mehrwert zu bieten, sollte der Fokus im nächsten Schritt aber unbedingt darauf liegen, die Daten nativ in der App darzustellen, anstatt auf externe Viewer zu verweisen.